\section{Seguridad y alineación}

\subsection{La IA necesita ser «mostrada», no solo discutida}

El pensamiento de Ilya cambió de manera importante durante el último año: ahora le da mucha más importancia al \textbf{despliegue incremental} de la IA. La razón principal es que a las personas (incluidos quienes trabajan en IA) les resulta muy difícil imaginar cómo será la IA del futuro.

\begin{importantbox}{Mostrar la IA frente a discutir la IA}
"Suppose you read an essay about AI... Now suppose you \textit{see} an AI doing this and that. It is incomparable."

Leer un artículo sobre IA frente a ver con los propios ojos las capacidades de la IA: ambas cosas son completamente incomparables. Por eso la IA necesita desplegarse frente al público, en lugar de saltar directamente a la superinteligencia dentro del laboratorio.
\end{importantbox}

\subsection{Tres predicciones clave}

\begin{enumerate}
    \item \textbf{Los competidores empezarán a cooperar en materia de seguridad}: ya hay indicios de esto (la cooperación inicial entre OpenAI y Anthropic)
    \item \textbf{Cuando la IA empiece a «feel powerful», todas las empresas se volverán más paranoicas}: en la actualidad la IA comete errores, por lo que no se percibe como poderosa, pero esto va a cambiar
    \item \textbf{Los gobiernos y el público exigirán que se tomen medidas}: a medida que el poder de la IA se vuelva más evidente
\end{enumerate}

\subsection{Cuidar de los seres sintientes (Sentient Life)}

\begin{importantbox}{Objetivo de alineación: cuidar de los seres sintientes}
Ilya propone un objetivo de alineación concreto: hacer que la IA \textbf{cuide de los seres sintientes} (care about sentient life), y no solo de los seres humanos. Las razones son:
\begin{itemize}
    \item Es posible que la IA misma sea sintiente (sentient), y que sea más fácil lograr que cuide de «los seres sintientes, incluida ella misma» que lograr que cuide únicamente de los seres humanos
    \item Por analogía con la empatía humana hacia los animales, que quizás proviene de que «usamos los mismos circuitos neuronales para simular tanto a nosotros mismos como a los demás»
    \item Si los primeros N sistemas de superinteligencia realmente cuidan de los seres sintientes, entonces las cosas podrían funcionar bien durante bastante tiempo, al menos
\end{itemize}
\end{importantbox}

\begin{warningbox}{El dilema del equilibrio a largo plazo}
Ilya reconoce que el equilibrio a largo plazo es difícil. Un escenario posible: cada persona tendría una IA que ganaría dinero por ella y la representaría en el ámbito político, pero los seres humanos dejarían de ser participantes reales. Él propone una solución que «a él mismo no le gusta pero que hay que considerar»: \textbf{que los seres humanos se conviertan en «medio IA» mediante algún tipo de neural link++}, de modo que la comprensión de la IA se transmita directamente a los seres humanos, manteniendo así su participación dentro del ciclo.
\end{warningbox}

\subsection{Limitar el poder de la superinteligencia más fuerte}

\begin{importantbox}{Cap the Power}
"I think it would be really materially helpful if the power of the most powerful super intelligence was somehow capped."

Ilya considera que limitar de alguna forma la capacidad del sistema más poderoso sería «materially helpful», aunque reconoce no estar seguro de cómo lograrlo. Señala analogías con la sociedad humana: los seres humanos son «medio agent» (persiguen una recompensa, se cansan de ella y cambian a otra); el mercado es un agent miope; la evolución es a la vez inteligente y torpe; los gobiernos están diseñados como un juego perpetuo entre tres poderes.
\end{importantbox}

\subsection{Resumen de la sección}

La visión de seguridad de Ilya evolucionó durante el último año: ahora da más importancia al despliegue incremental y a mostrar públicamente las capacidades de la IA. Propone «cuidar de los seres sintientes» como objetivo de alineación, reconoce que el equilibrio a largo plazo es difícil, y considera que limitar el poder de los sistemas más fuertes sería «materially helpful».

